% GENERATED FILE. Edit canonical lessons, then run npm run generate:books.
% canonical-source-hash: fnv1a64:a17b3b3d0224b969
% canonical-lessons: RU-C360-soveshchanie, RU-C360-otchyot, RU-C360-grafik, RU-C360-smena, RU-C360-otchityvatsya, RU-R360-review-353-360

\chapter{Meetings and reports}
\label{ch:ru-meetings-and-reports}
\hlchaptermodality{360}

\begin{chapteropening}
\textbf{By the end of this chapter:} \emph{I can talk about a work meeting, a report, my schedule and my shift, and say whom I report to.}

Complete the last lesson of chapter 360: I can talk about a work meeting, a report, my schedule and my shift, and say whom I report to.
\end{chapteropening}

\section[\texorpdfstring{\ru{совещание} (soveshchániye)}{совещание (soveshchániye)}]{\ru{совещание} (soveshchániye) — a meeting (at work)}

\label{lesson:RU-C360-soveshchanie}



\begin{quote}
Before the new one: say the Russian for to carry out, then the Russian for to cope.
\end{quote}

\subsection*{You'll want to know: \ru{совещание}}
\textbf{\ru{совещание}} (\emph{soveshchániye}) — \textquotedblleft{}a meeting (at work)\textquotedblright{}.

\begin{quote}
\ru{Совещание} \ru{начнётся} \ru{в} \ru{девять} \ru{часов}.
\end{quote}

\begin{quote}
\emph{Soveshchániye nachnyótsya v dévyat chasóv.}
\end{quote}

\begin{quote}
The meeting will begin at nine o’clock.
\end{quote}

Neuter. A meeting at work; a meeting of friends is \emph{\ru{встреча}}.

\begin{cousinweb}
\textbf{\ru{совещание}} is built on \emph{\ru{совещаться}} \textquotedblleft{}to confer\textquotedblright{}, on \emph{\ru{вещать}} \textquotedblleft{}to proclaim\textquotedblright{}, the Church Slavonic form of the root of \emph{\ru{ответ}}, Old Russian \emph{větŭ} \textquotedblleft{}counsel\textquotedblright{}. The \emph{\ru{щ}} is the Church Slavonic shape where Russian has \emph{\ru{ч}}.
\end{cousinweb}

\begin{grammarlens}[title={What you've built}]
One more word for telling, explaining and describing at B1.
\end{grammarlens}

\subsection*{Your turn}
\begin{itemize}
\raggedright
  \item \emph{Say it:} \emph{soveshchániye}
  \item \emph{Say it:} \emph{soveshchániye}, once more
  \item \emph{Say it:} a sentence of your own with \emph{soveshchániye}
  \item \emph{Recall:} say the Russian for to carry out, then the Russian for to cope, then say \emph{soveshchániye} again
\end{itemize}

\begin{tcolorbox}[breakable,colback=teal!4,colframe=teal!35!black,title={Before you move on}]
What does \ru{совещание} mean? (\textquotedblleft{}A meeting (at work)\textquotedblright{}.) Say it once more.
\end{tcolorbox}
\section[\texorpdfstring{\ru{отчёт} (otchyót)}{отчёт (otchyót)}]{\ru{отчёт} (otchyót) — a report}

\label{lesson:RU-C360-otchyot}



\begin{quote}
Before the new one: say the Russian for to cope, then the Russian for a meeting.
\end{quote}

\subsection*{You'll want to know: \ru{отчёт}}
\textbf{\ru{отчёт}} (\emph{otchyót}) — \textquotedblleft{}a report\textquotedblright{}.

\begin{quote}
\ru{Я} \ru{пишу} \ru{отчёт} \ru{о} \ru{командировке}.
\end{quote}

\begin{quote}
\emph{Ya pishú otchyót o komandiróvke.}
\end{quote}

\begin{quote}
I am writing a report on the business trip.
\end{quote}

Masculine. \emph{\ru{Сдать} \ru{отчёт}} — to hand in a report. \emph{\ru{Отдавать} \ru{себе} \ru{отчёт}} — to realise.

\begin{cousinweb}
\textbf{\ru{отчёт}} is built from \emph{\ru{от}-} and the root of \emph{\ru{счёт}} \textquotedblleft{}a count, an account\textquotedblright{} and \emph{\ru{считать}} \textquotedblleft{}to count\textquotedblright{}: a report is an account given back.
\end{cousinweb}

\begin{grammarlens}[title={What you've built}]
One more word for telling, explaining and describing at B1.
\end{grammarlens}

\subsection*{Your turn}
\begin{itemize}
\raggedright
  \item \emph{Say it:} \emph{otchyót}
  \item \emph{Say it:} \emph{otchyót}, once more
  \item \emph{Say it:} a sentence of your own with \emph{otchyót}
  \item \emph{Recall:} say the Russian for to cope, then the Russian for a meeting, then say \emph{otchyót} again
\end{itemize}

\begin{tcolorbox}[breakable,colback=teal!4,colframe=teal!35!black,title={Before you move on}]
What does \ru{отчёт} mean? (\textquotedblleft{}A report\textquotedblright{}.) Say it once more.
\end{tcolorbox}
\section[\texorpdfstring{\ru{график} (gráfik)}{график (gráfik)}]{\ru{график} (gráfik) — a schedule; a graph}

\label{lesson:RU-C360-grafik}



\begin{quote}
Before the new one: say the Russian for a meeting, then the Russian for a report.
\end{quote}

\subsection*{You'll want to know: \ru{график}}
\textbf{\ru{график}} (\emph{gráfik}) — \textquotedblleft{}a schedule; a graph\textquotedblright{}.

\begin{quote}
\ru{У} \ru{меня} \ru{свободный} \ru{график} \ru{работы}.
\end{quote}

\begin{quote}
\emph{U menyá svobódnyy gráfik rabóty.}
\end{quote}

\begin{quote}
I have flexible working hours.
\end{quote}

Masculine. \emph{\ru{Гибкий} \ru{график}} — flexible hours. It is also a chart of figures.

\begin{cousinweb}
\textbf{\ru{график}} is usually traced, through German, to Greek \emph{graphikós} \textquotedblleft{}of drawing or writing\textquotedblright{}, built on \emph{gráphein} \textquotedblleft{}to scratch, to draw, to write\textquotedblright{}. A schedule is a time plan drawn out.
\end{cousinweb}

\begin{grammarlens}[title={What you've built}]
One more word for telling, explaining and describing at B1.
\end{grammarlens}

\subsection*{Your turn}
\begin{itemize}
\raggedright
  \item \emph{Say it:} \emph{gráfik}
  \item \emph{Say it:} \emph{gráfik}, once more
  \item \emph{Say it:} a sentence of your own with \emph{gráfik}
  \item \emph{Recall:} say the Russian for a meeting, then the Russian for a report, then say \emph{gráfik} again
\end{itemize}

\begin{tcolorbox}[breakable,colback=teal!4,colframe=teal!35!black,title={Before you move on}]
What does \ru{график} mean? (\textquotedblleft{}A schedule; a graph\textquotedblright{}.) Say it once more.
\end{tcolorbox}
\section[\texorpdfstring{\ru{смена} (sména)}{смена (sména)}]{\ru{смена} (sména) — a shift (at work); a change}

\label{lesson:RU-C360-smena}



\begin{quote}
Before the new one: say the Russian for a report, then the Russian for a schedule.
\end{quote}

\subsection*{You'll want to know: \ru{смена}}
\textbf{\ru{смена}} (\emph{sména}) — \textquotedblleft{}a shift (at work); a change\textquotedblright{}.

\begin{quote}
\ru{Сегодня} \ru{у} \ru{медсестры} \ru{ночная} \ru{смена}.
\end{quote}

\begin{quote}
\emph{Segódnya u medsestrý nochnáya sména.}
\end{quote}

\begin{quote}
Today the nurse has a night shift.
\end{quote}

Feminine. \emph{\ru{Работать} \ru{в} \ru{две} \ru{смены}} — to work in two shifts. \emph{\ru{Смена} \ru{погоды}} — a change in the weather.

\begin{cousinweb}
\textbf{\ru{смена}} is built on \emph{\ru{менять}} \textquotedblleft{}to change\textquotedblright{}, inherited from Proto-Slavic \emph{*měniti}: a shift is one team changing places with another.
\end{cousinweb}

\begin{grammarlens}[title={What you've built}]
One more word for telling, explaining and describing at B1.
\end{grammarlens}

\subsection*{Your turn}
\begin{itemize}
\raggedright
  \item \emph{Say it:} \emph{sména}
  \item \emph{Say it:} \emph{sména}, once more
  \item \emph{Say it:} a sentence of your own with \emph{sména}
  \item \emph{Recall:} say the Russian for a report, then the Russian for a schedule, then say \emph{sména} again
\end{itemize}

\begin{tcolorbox}[breakable,colback=teal!4,colframe=teal!35!black,title={Before you move on}]
What does \ru{смена} mean? (\textquotedblleft{}A shift (at work); a change\textquotedblright{}.) Say it once more.
\end{tcolorbox}
\section[\texorpdfstring{\ru{отчитываться} (otchítyvatsya)}{отчитываться (otchítyvatsya)}]{\ru{отчитываться} (otchítyvatsya) — to report, to account for}

\label{lesson:RU-C360-otchityvatsya}



\begin{quote}
Before the new one: say the Russian for a schedule, then the Russian for a shift.
\end{quote}

\subsection*{You'll want to know: \ru{отчитываться}}
\textbf{\ru{отчитываться}} (\emph{otchítyvatsya}) — \textquotedblleft{}to report, to account for\textquotedblright{}.

\begin{quote}
\ru{Каждую} \ru{пятницу} \ru{мы} \ru{отчитываемся} \ru{перед} \ru{начальником}.
\end{quote}

\begin{quote}
\emph{Kázhduyu pyátnitsu my otchítyvayemsya péred nachálnikom.}
\end{quote}

\begin{quote}
Every Friday we report to the boss.
\end{quote}

Imperfective; perfective \emph{\ru{отчитаться}}. \emph{\ru{Отчитываться} \ru{перед} \ru{кем} \ru{о} \ru{чём}}. Without \emph{-\ru{ся}}, \emph{\ru{отчитывать}} means to tell off.

\begin{cousinweb}
\textbf{\ru{отчитываться}} is built from \emph{\ru{от}-} and \emph{\ru{читать}} in its old sense \textquotedblleft{}to count\textquotedblright{}, the root of \emph{\ru{число}} \textquotedblleft{}number\textquotedblright{}, Proto-Slavic \emph{*čislo} (EtymDB): to give back an account.
\end{cousinweb}

\begin{grammarlens}[title={What you've built}]
One more word for telling, explaining and describing at B1.
\end{grammarlens}

\subsection*{Your turn}
\begin{itemize}
\raggedright
  \item \emph{Say it:} \emph{otchítyvatsya}
  \item \emph{Say it:} \emph{otchítyvatsya}, once more
  \item \emph{Say it:} a sentence of your own with \emph{otchítyvatsya}
  \item \emph{Recall:} say the Russian for a schedule, then the Russian for a shift, then say \emph{otchítyvatsya} again
\end{itemize}

\begin{tcolorbox}[breakable,colback=teal!4,colframe=teal!35!black,title={Before you move on}]
What does \ru{отчитываться} mean? (\textquotedblleft{}To report, to account for\textquotedblright{}.) Say it once more.
\end{tcolorbox}
\section[(dialogue)]{Words from chapters 353–360 — a review}

\label{lesson:RU-R360-review-353-360}



\begin{quote}
Say the Russian for damage and to cause.
\end{quote}

\subsection*{Your turn}
Say these aloud:

\begin{itemize}
\raggedright
  \item damage, to cause, to harm, to suffer, to make worse
  \item to improve, to increase, to decrease, to lower, to raise
  \item accidental, chance, a coincidence, inevitable, to arise
  \item despite, nevertheless, anyway, on the one hand, on the other hand
  \item an employee, a manager, a subordinate, to manage, a team
  \item a vacancy, a job interview, a candidate, to hire, length of service
  \item a duty, to entrust, an assignment, to carry out, to cope
  \item a meeting, a report, a schedule, a shift, to report
\end{itemize}

\begin{tcolorbox}[breakable,colback=teal!4,colframe=teal!35!black,title={Before you move on}]
Damage? (\textbf{\ru{ущерб}}.) To cause? (\textbf{\ru{причинять}}.) To harm? (\textbf{\ru{вредить}}.) To suffer? (\textbf{\ru{страдать}}.) To make worse? (\textbf{\ru{ухудшать}}.) To improve? (\textbf{\ru{улучшать}}.) To increase? (\textbf{\ru{увеличиваться}}.) To decrease? (\textbf{\ru{уменьшаться}}.) To lower? (\textbf{\ru{снижать}}.) To raise? (\textbf{\ru{повышать}}.) Accidental? (\textbf{\ru{случайный}}.) Chance? (\textbf{\ru{случайность}}.) A coincidence? (\textbf{\ru{совпадение}}.) Inevitable? (\textbf{\ru{неизбежный}}.) To arise? (\textbf{\ru{возникать}}.) Despite? (\textbf{\ru{несмотря} \ru{на}}.) Nevertheless? (\textbf{\ru{тем} \ru{не} \ru{менее}}.) Anyway? (\textbf{\ru{впрочем}}.) On the one hand? (\textbf{\ru{с} \ru{одной} \ru{стороны}}.) On the other hand? (\textbf{\ru{с} \ru{другой} \ru{стороны}}.) An employee? (\textbf{\ru{сотрудник}}.) A manager? (\textbf{\ru{руководитель}}.) A subordinate? (\textbf{\ru{подчинённый}}.) To manage? (\textbf{\ru{руководить}}.) A team? (\textbf{\ru{коллектив}}.) A vacancy? (\textbf{\ru{вакансия}}.) A job interview? (\textbf{\ru{собеседование}}.) A candidate? (\textbf{\ru{кандидат}}.) To hire? (\textbf{\ru{нанимать}}.) Length of service? (\textbf{\ru{стаж}}.) A duty? (\textbf{\ru{обязанность}}.) To entrust? (\textbf{\ru{поручать}}.) An assignment? (\textbf{\ru{поручение}}.) To carry out? (\textbf{\ru{выполнять}}.) To cope? (\textbf{\ru{справляться}}.) A meeting? (\textbf{\ru{совещание}}.) A report? (\textbf{\ru{отчёт}}.) A schedule? (\textbf{\ru{график}}.) A shift? (\textbf{\ru{смена}}.) To report? (\textbf{\ru{отчитываться}}.)
\end{tcolorbox}
